\documentclass[10pt,a4paper]{article}
\usepackage[utf8]{inputenc}
\usepackage[a4paper,left=3cm,right=2cm]{geometry}
\usepackage{amsmath}
\usepackage[thmmarks, amsmath, thref]{ntheorem}
\usepackage{amsfonts}
\usepackage{amssymb}
\usepackage{fancyhdr}
\pagestyle{fancy}
\usepackage{anysize}
\usepackage{graphicx}
\usepackage{float}
\usepackage{hyperref}
\usepackage{enumitem}
\usepackage{verbatim}
\usepackage{siunitx}
\newtheorem{theorem}{Theorem}
\theoremsymbol{\ensuremath{\blacksquare}}
\newtheorem*{solution}{Solución:}

\usepackage{xcolor}

% Margenes y formalidades

\textwidth 6.25in % Margen Derecho y ancho del texto
\textheight 8.5in
\oddsidemargin 0in % Margen Izquierdo
\headheight 0in % Largo del texto
\leftmargin 1in
\parindent 0em % Salto entre la indentación (Sangría)https://www.overleaf.com/4519127669sdpptkhvsbfn
\topmargin -0.5in % Margen Superior
\parskip 0.5ex % Salto entre parrafos
\baselineskip 1pt



\lhead{
\includegraphics[scale=0.08]{figuras/logoUSM.pn}} % texto izquierda de la cabecera
%\chead{}                                      % texto centro de la cabecera
\rhead{Universidad Técnica Federico Santa María\\Ingeniería Comercial } % número de página a la derecha
\renewcommand{\headrulewidth}{0.4pt} % grosor de la línea de la cabecera
\renewcommand{\footrulewidth}{0.4pt} % grosor de la línea del pie

\begin{document}
\vspace{5cm}



\begin{center}
\begin{large}
\textbf{\\Control 3}
\end{large}\\
\end{center}

\begin{center}
\begin{large}
\textbf{Econometría ICS-294}
\end{large}\\
\vspace{0.2cm}
III Trimestre 2025\\
Profesor: Pedro Comte\\
%Ayudante: 
\end{center}

\textbf{Instrucciones:} 

\begin{itemize}

    \item Usted dispone de 60 minutos  para intentar los 100 puntos del control.
    \item Cuide la presentación y redacción de sus preguntas.
    \item No puede utilizar celular, computador o tablet, pero sí una calculadora TI.
     
\end{itemize}

\section{Preguntas cortas (40 puntos)}

\begin{enumerate}



\item (20 puntos) 
Se presenta el siguiente modelo econométrico para explicar las importaciones (IMP) en función del Producto Interno Bruto (PIB), la Inversión (INV) y el Consumo (CONS):

\[
\text{IMP} = -19.7251 + 0.0322 \cdot \text{PIB} + 0.4142 \cdot \text{INV} + 0.2427 \cdot \text{CONS}
\]

Considerando la información disponible, responda lo siguiente:

\begin{enumerate}
    \item (10 puntos) ¿Existe evidencia de multicolinealidad entre las variables explicativas? Fundamente su respuesta.

Para su respuesta considere la siguiente información.

    Matriz de correlaciones:

    \[
    \begin{array}{c|ccc}
    & \text{PIB} & \text{INV} & \text{CONS} \\
    \hline
    \text{PIB} & 1 & 0.2154 & 0.9989 \\
    \text{INV} & 0.2154 & 1 & 0.2137 \\
    \text{CONS} & 0.9989 & 0.2137 & 1 \\
    \end{array}
    \]

    Coeficientes de determinación de regresiones auxiliares para cada regresor:

    \[
    R_{PIB}^2 = 0.9979, \quad R_{INV}^2 = 0.0475, \quad R_{CONS}^2 = 0.9979
    \]

{\color{blue}
\textbf{Sí, existe evidencia clara de multicolinealidad, particularmente entre las variables \textit{PIB} y \textit{CONS}.} Esta conclusión se fundamenta en (BASTA QUE EXPLIQUEN 1):

\begin{itemize}
    \item \textbf{Matriz de correlaciones}

    La correlación entre \textit{PIB} y \textit{CONS} es $0{,}9989$, lo que indica una relación casi perfecta entre ambas variables.

    Las correlaciones entre \textit{INV} y \textit{PIB} ($0{,}2154$) e \textit{INV} y \textit{CONS} ($0{,}2137$) son bajas, lo que sugiere que \textit{INV} no presenta multicolinealidad severa.

    \item \textbf{Coeficientes de determinación de las regresiones auxiliares}

    El coeficiente de determinación $R^2$ de las regresiones auxiliares mide cuán bien una variable explicativa puede ser explicada por las otras:

    \[
    R^2_{\text{PIB}} = 0{,}9979, \quad R^2_{\text{CONS}} = 0{,}9979, \quad R^2_{\text{INV}} = 0{,}0475
    \]

    Un valor de $R^2$ cercano a 1 en las regresiones auxiliares de \textit{PIB} y \textit{CONS} indica que ambas variables están casi completamente explicadas por las otras variables del modelo.

    En particular, el alto $R^2$ sugiere que \textit{PIB} y \textit{CONS} están prácticamente colineales, lo que puede causar inestabilidad en las estimaciones del modelo.

    \item \textbf{Criterio del Factor de Inflación de la Varianza (VIF)} 

    Recordando que:

    \[
    \text{VIF} = \frac{1}{1 - R_j^2}
    \]

    Para \textit{PIB} y \textit{CONS}:

    \[
    \text{VIF} \approx \frac{1}{1 - 0{,}9979} \approx \frac{1}{0{,}0021} \approx 476{,}2
    \]

    Un VIF mayor a 10 ya es señal de multicolinealidad severa. En este caso, el VIF está sobre 400, lo cual confirma una multicolinealidad extremadamente alta.
\end{itemize}

{\color{black}
    \item (10 puntos) En caso de confirmarse la existencia de multicolinealidad, ¿qué métodos o estrategias podrían emplearse para mitigar este problema en el modelo?
\end{enumerate}

(BASTA QUE EXPLIQUEN 1)
{\color{blue}
\begin{itemize}
    \item \textbf{Eliminar una de las variables altamente correlacionadas}
    \begin{itemize}
        \item Dado que \textit{PIB} y \textit{CONS} están altamente correlacionadas, se puede considerar eliminar una de las dos (por ejemplo, \textit{CONS} si está incluido implícitamente en el \textit{PIB}).
        \item Esto puede mejorar la precisión de las estimaciones sin perder poder explicativo.
    \end{itemize}

    \item \textbf{Combinar variables en un índice}
    \begin{itemize}
        \item Se puede construir un índice compuesto que resuma la información común entre \textit{PIB} y \textit{CONS}.
        \item Esto reduce la colinealidad al incorporar solo una nueva variable no correlacionada con las demás.
    \end{itemize}

    \item \textbf{Reespecificación del modelo}
    \begin{itemize}
        \item Puede implicar usar tasas de crecimiento o transformaciones (como logaritmos) que reduzcan la correlación entre variables.
    \end{itemize}

    \item \textbf{Recolección de más datos}
    \begin{itemize}
        \item Si es posible, aumentar el tamaño de la muestra puede ayudar a mejorar la precisión de las estimaciones, aunque no elimina la multicolinealidad.
    \end{itemize}
\end{itemize}
{\color{black}
\item (20 puntos)
¿Cómo podría utilizar el test de White (simplificado) para detectar heterocedasticidad en el modelo \( y_i = \beta_0 + \beta_1 x_{1i} + \beta_2 x_{2i} + u_i \)? Describa paso a paso el procedimiento e indique cómo se interpreta el resultado

{\color{blue}
\begin{enumerate}
    \item \textbf{Estimar el modelo original por MCO}

    \[
    y_i = \hat{\beta}_0 + \hat{\beta}_1 x_{1i} + \hat{\beta}_2 x_{2i} + \hat{u}_i
    \]

    Obtienes los residuos estimados \( \hat{u}_i \).

    \item \textbf{Elevar al cuadrado los residuos}

    Calcula:

    \[
    \hat{u}_i^2
    \]

    Esta será la variable dependiente en la regresión auxiliar.

    \item \textbf{Regresión auxiliar simplificada}

    Estima la siguiente regresión:

    \[
    \hat{u}_i^2 = \alpha_0 + \alpha_1 \hat{y}_{i}^2 + \alpha_2 \hat{y}^2_{i} + e_i
    \]

    \textit{Nota:} No se incluyen los términos originales ni los productos cruzados, sólo los cuadrados de los valores predichos de y.

    \item \textbf{Obtener el estadístico LM}

    Calcula el estadístico:

    \[
    LM = n \cdot R^2_{\text{aux}}\sim X^2_{df=2}
    \]

    Donde:

    \begin{itemize}
        \item \( n \): número de observaciones,
        \item \( R^2_{\text{aux}} \): coeficiente de determinación de la regresión auxiliar.

         El estadístico \( LM \) sigue una distribución \( \chi^2 \) con 2 grados de libertad.
    \end{itemize}

    \item \textbf{Hipótesis y decisión}

\[
H_0: \alpha_1=\alpha_2 = 0
\]
      
\[
 \alpha_1\neq 0 \quad  o \quad\alpha_2 \neq 0
\]


   

    \textbf{Decisión:}
    
    \begin{itemize}
        \item Si \( LM > \chi^2_{\alpha, 2} \), se \textbf{rechaza} \( H_0 \) → hay heterocedasticidad.
        \item Si \( LM \leq \chi^2_{\alpha, 2} \), \textbf{no se rechaza} \( H_0 \) → no hay evidencia de heterocedasticidad.
    \end{itemize}
\end{enumerate}

\end{enumerate}
\color{black}
\section{RdA3.2  Ejercicios (60 puntos)}


\begin{enumerate}

\item (30 puntos) A continuación se presenta el gráfico de la serie temporal del tipo de cambio nominal (pesos por dólar) en Chile, correspondiente al período comprendido entre enero de 1983 y diciembre de 2024, con datos de frecuencia mensual. 


\begin{center}
\textbf{Figura 1: Tipos de cambio (pesos por dólar)}
\begin{figure}[H]
    \centering
    \includegraphics[width=0.85\linewidth]{certamenes/BCChChart.png}
\end{figure}
\end{center}

Con el objetivo de evaluar si la serie temporal sigue un proceso de caminata aleatoria (random walk), un grupo de investigadores aplica el test de Dickey-Fuller mediante la estimación de dos modelos auxiliares:

   \textbf{Modelo sin tendencia (without trend):}
    \begin{table}[H]
    \centering
    \begin{tabular}{lccc}
    \toprule
    Variable & Coef. & Error estándar & Estadístico t \\
    \midrule
    Intercepto & 0,0710 & 0,0354 & 2,006 \\
    $y_{t-1}$  & -0,0182 & 0,0097 & -1,876 \\
    \bottomrule
    \end{tabular}
    \end{table}

    \textbf{Modelo auxiliar con tendencia (\textit{trend}):}
    \begin{table}[H]
    \centering
    \begin{tabular}{lccc}
    \toprule
    Variable & Coef. & Error estándar & Estadístico t \\
    \midrule
    Intercepto & 0,0452 & 0,0289 & 1,563 \\
    Tendencia ($t$) & 0,0021 & 0,0003 & 7,000 \\
    $y_{t-1}$  & -0,0304 & 0,018 & -1,6889 \\
    \bottomrule
    \end{tabular}
    \end{table}

  Testee con un 1\% de significancia si el tipo cambio nominal tiene raíz unitaria. Explique su resultado.  


{\color{blue}
  
La inspección visual de la serie permite identificar a un crecimiento sistemático a lo largo del tiempo, compatible con una tendencia temporal. Por tanto, se justifica el uso del test Dickey-Fuller con tendencia.


Modelo auxiliar: \\
\[ 
\Delta y_t = \alpha + \beta t + \phi y_{t-1} + \varepsilon_t 
\]
Hipótesis:
\begin{itemize}
    \item \( H_0: \phi = 0 \) → la serie tiene raíz unitaria (no estacionaria).
    \item \( H_1: \phi < 0 \) → la serie es estacionaria alrededor de una tendencia determinista.
\end{itemize}


Si se rechaza \( H_0 \), se concluye que el proceso es estacionario alrededor de una tendencia determinista, es decir, fluctúa de forma estable en torno a una línea recta en el tiempo.


\begin{itemize}
    \item Sin tendencia: Estadístico t para \( y_{t-1} \) = $-1{,}876$
    \item Con tendencia: Estadístico t para \( y_{t-1} \) = $-1{,}6889$
\end{itemize}


Se utiliza el valor del estadístico t del modelo con tendencia: $-1{,}6889$.\\

La serie tiene una frecuencia mensual desde enero de 1983 a diciembre de 2024  con un total de 504 observaciones.

Se deben utilizar los valores críticos de la tabla para T=500:

\begin{itemize}
    \item 1\%: $-3{,}98$
    \item 5\%: $-3{,}42$
\end{itemize}

Comparación: \(-1{,}6889 > -3{,}98\) → No se rechaza \( H_0 \)  al 1\%.\\

La evidencia no permite rechazar la presencia de raíz unitaria. La serie es no estacionaria y sigue una caminata aleatoria con tendencia estocástica.
}


\item (30 puntos) Suponga que usted está interesado en determinar los factores que explican
las diferencias salariales entre distintos individuos. Para realizar
su análisis, usted especifica la siguiente función de regresión:

\begin{equation}
wage = \beta_{1} + \beta_{2}IQ + \beta_{3}Educ + \beta_{4}Exper + \beta_{5}Tenure + \beta_{6}personacolor + u \label{eq:original}
\end{equation}

donde \emph{wage} representa el salario, \emph{IQ} es el coeficiente intelectual,
\emph{Educ} son los años de educación, \emph{Exper} los años de experiencia laboral,
\emph{Tenure} la antigüedad en el trabajo actual, y \emph{personacolor} es una variable dicotómica
que toma el valor 1 cuando la persona se identifica como {persona de color}, y 0 en caso contrario.


Los resultados de la regresión se presentan en la Tabla 1, y en la Figura 1 se muestran los residuos
agrupados según la variable \emph{personacolor}:

\begin{center}
\textbf{Tabla 1}
\end{center}

\begin{center}
Model 1: OLS, using observations 1--935\\ Dependent variable: wage\\
\vspace{1em}
\begin{tabular}{lr@{.}lr@{.}lr@{.}lr@{.}l}
 &  \multicolumn{2}{c}{Coefficient} &   \multicolumn{2}{c}{Std.\ Error} &    \multicolumn{2}{c}{$t$-ratio} &     \multicolumn{2}{c}{p-value} \\
\hline
const &   $-$405&803 &     122&400 &       $-$3&3154 &         0&0010 \\
IQ &   3&77957 &     0&996294 &       3&7936 &         0&0002 \\
educ &   57&5438 &     7&00006 &       8&2205 &         0&0000 \\
exper &   14&8737 &     3&19330 &       4&6578 &         0&0000 \\
tenure &   6&88539 &     2&46227 &       2&7964 &         0&0053 \\
personacolor & -40&15 & 38&70 & -1&037 & 0&296\\
\end{tabular}
\vspace{1ex}

\begin{tabular}{lrlr}
Mean dependent var &  957.9455 & S.D. dependent var &  404.3608 \\
Sum squared resid &  1.25\textrm{e+08} & S.E. of regression &  367.4213 \\
$R^2$ &  0.178780 & Adjusted $R^2$ &  0.174361 \\
$F(5, 929)$ &  40.44887 & P-value($F$) &  1.07\textrm{e--37} \\
\end{tabular}
\end{center}

\begin{center}
\textbf{Figura 1: Residuos agrupados por identificación étnico-racial}
\begin{figure}[H]
    \centering
    \includegraphics[width=0.65\linewidth]{certamenes/racial.png}
\end{figure}
\end{center}

\begin{enumerate}
    \item (6 puntos) De acuerdo a esta información, ¿le parece que existe heterocedasticidad en el modelo? ¿Tiene sentido este resultado? Explique.

 

    \item (10 puntos)  En la Tabla 2 se presentan los resultados de una regresión de los residuos al cuadrado de la regresión de la tabla 1 sobre las variables explicativas:

\begin{center}
\textbf{Tabla 2: OLS, usando 935 observaciones} \\
Variable dependiente: $\hat{u}^2$ \\
\vspace{1em}
\begin{tabular}{lr@{.}lr@{.}lr@{.}lr@{.}l}
 & \multicolumn{2}{c}{Coeficiente} & \multicolumn{2}{c}{Error estándar} & \multicolumn{2}{c}{$t$-ratio} & \multicolumn{2}{c}{p-valor} \\[1ex]
const          & $-$160277&000 & 104551&00 & $-$1&5330 & 0&1256 \\
IQ             & 921&598   & 851&010  & 1&0829  & 0&2791 \\
educ           & 16288&600 & 5979&280 & 2&7242  & 0&0066 \\
exper          & 2955&490  & 2727&640 & 1&0835  & 0&2789 \\
tenure         & $-$6586&500 & 2103&220 & $-$3&1316 & 0&0018 \\
personacolor   & $-$70000&000 & 30000&000 & $-$2&3333 & 0&0200 \\
\end{tabular}

\vspace{1ex}
\begin{tabular}{lrlr}
Media variable dependiente & 134132.1 & Desv. estándar dependiente & 317721.3 \\
Suma de residuos al cuadrado & $9.15 \times 10^{13}$ & Error estándar de regresión & 313842.4 \\
$R^2$ & 0.029491 & $R^2$ ajustado & 0.024268 \\
$F(5, 929)$ & 5.646001 & Valor-p ($F$) & 0.000039 \\
\end{tabular}
\end{center}


    \[
    \hat{\mu}^2 = a_1 + a_2 IQ + a_3 Educ + a_4 Exper + a_5 Tenure + a_6 personacolor + v
    \]

    ¿Puede usarse esta regresión para testear la presencia de heterocedasticidad? Explique e interprete.

      \item (8 puntos) Observe que la variable \texttt{personacolor} no es estadísticamente significativa según los resultados de la tabla 1. ¿Podemos concluir que \texttt{personacolor} no afecta el ingreso? Fundamente su respuesta considerando la presencia de heterocedasticidad. 

        \item (6 puntos) ¿Cómo se puede corregir el problema anterior para obtener inferencias válidas? ¿Qué tipo de errores estándar se deben utilizar? 

   
\end{enumerate}

\end{enumerate}

\begin{enumerate}
    \item {\color{blue}
    A partir de la Figura 1, se observa que la dispersión de los residuos parece variar dependiendo del valor de la variable \emph{personacolor}. En particular, los residuos correspondientes a personas de color presentan una mayor dispersión en comparación con el resto de la muestra, lo cual sugiere la presencia de heterocedasticidad.

    Este resultado tiene sentido desde el punto de vista económico si se considera que las personas de color pueden enfrentar condiciones laborales más heterogéneas, lo que se traduce en una mayor variabilidad no explicada del salario.}

    \item {\color{blue}
    Esta regresión puede utilizarse para testear la presencia de heterocedasticidad utilizando el test de White. La idea central es estimar una regresión auxiliar donde los residuos al cuadrado $\hat{u}^2$ se explican en función de las variables originales del modelo.

    En este caso, se aplica el test de White a través del siguiente contraste:

    \[
    H_0: a_2 = a_3 = a_4 = a_5 = a_6 = 0 \quad \text{(homocedasticidad)}
    \]
    \[
    H_1: \text{al menos un } a_j \neq 0 \quad \text{(heterocedasticidad)}
    \]

    Se utiliza como estadístico de contraste el multiplicador de Lagrange:

    \[
    LM = n \cdot R^2 = 935 \cdot 0.029491 \approx 27.57
    \]

    Este estadístico se distribuye asintóticamente como una $\chi^2$ con 5 grados de libertad (tantos como regresores excluyendo la constante). El valor crítico al 5\% es aproximadamente 11.07.

    Dado que $27.57 > 11.07$, se rechaza la hipótesis nula de homocedasticidad. Por lo tanto, se concluye que existe evidencia estadísticamente significativa de heterocedasticidad en el modelo.}


 \item {\color{blue}
    No necesariamente. Dado que hay evidencia de heterocedasticidad, el valor-p que indica que \texttt{color} no es significativa puede ser inválido. La heterocedasticidad puede haber inflado o reducido artificialmente el error estándar de ese coeficiente. Por tanto, no es posible concluir con certeza que \texttt{color} no tiene efecto.
    }   

  \item   {\color{blue}
    Para obtener inferencias válidas en presencia de heterocedasticidad, se deben utilizar errores estándar robustos. En particular, se recomienda emplear los errores estándar de tipo White (robustos a heterocedasticidad), que corrigen la matriz de varianzas-covarianzas sin asumir homocedasticidad.
    }
\end{enumerate}


\end{enumerate}








\newpage


\begin{center}
    {\Large Formulario  }

 
\end{center}

\begin{itemize}

% \item  
% {\color{red}
% $$\hat{\beta}_0 = \bar{y} - \hat{\beta}_1 \bar{x} $$
% $$\hat{\beta}_1 = \frac{\sum_{i=1}^{n}(x_i - \bar{x})(y_i-\bar{y})}{\sum_{i=1}^{n}(x_i - \bar{x})^2}$$}

\item  {\color{red} \[ \hat{\beta} = (X^{\top}X)^{-1}X^{\top}y \]}

\item 
	{\color{red}	\[ SST = \sum_{i=1}^{n} (y_i - \bar{y})^2 \quad SSE = \sum_{i=1}^{n} (\hat{y}_i - \bar{y})^2  \quad SSR = \sum_{i=1}^{n} \hat{u}_i^2\]}
		

    \item {\color{red} \[ R^2  = 1 - \frac{SSR}{SST} \quad \quad \bar{R}^2 =  1 - \frac{(n-1)SSR}{(n-k-1)SST} \]}

 \item {\color{red} \[ t_{\hat{\beta_j}}=\frac{\hat{\beta}_j - \beta_j}{\text{s.e.}(\hat{\beta}_j)} \sim t_{n-k-1} \]}

 \item {\color{red}\[
    F = \frac{SSR_{r} - SSR_{ur}}{SSR_{ur}} \cdot \frac{n - k - 1}{j} \sim F(j, n - k - 1)
    \]}

\item {\color{red}Test F para regresión auxiliar.}
\begin{center}
 {  \color{red}     \(F=\dfrac{R^{2}/j}{(1-R^{2})/(n-k-1)}
          \sim F(j,\;n-k-1)\).}
\end{center}
% \item {\color{red}Considere que el valor crítico a un nivel de significancia del 5\% es 1.645}
% \item {\color{red}Considere los siguientes valores críticos para t: $\alpha = 0.01 \implies t \approx 2.62$; $\alpha = 0.05 \implies t \approx 1,97$;$\alpha = 0.01 \implies t \approx 2.62$; $\alpha = 0.1 \implies t \approx 1,645$}
%\item {\color{red}Considere que el valor crítico de F al 1\% con grados de libertad (2,178) es aproximadamente 4.72}


\item \color{red}{ Multiplicador de Lagrange

$$LM = nR^2 \sim \chi^2_{df}$$

Tabla $\chi^2_{df}$
\begin{center}
    \begin{tabular}{|c|c|c|c|}
    \hline
        df & 0.99 & 0.95 & 0.90 \\ 
        \hline
        1 &  6.635 & 3.841 & 2.706\\
        \hline
        2 &  9.210 & 5.991 & 4.605 \\
        \hline
        3 & 11.345 & 7.815 & 6.251  \\
        \hline
        4 &  13.277 & 9.488 & 7.779 \\
        \hline
        5 &  15.086 & 11.070 & 9.236 \\
        \hline
        6 &  16.812 & 12.592 & 10645 \\
        \hline
    \end{tabular}
    \end{center}}


\item Tabla de test D-F
\begin{center}
    \begin{figure}[H]
        \centering
        \includegraphics[width=0.7\linewidth]{sem14/tabla.png}
    \end{figure}
\end{center}

\end{itemize}


% \begin{figure}[H]
%     \centering
%     \includegraphics[width=0.85\linewidth]{Ejercicios//ejercicios_inferencia/cert2_F2_1184.png}
% \end{figure}

% \begin{figure}[H]
%     \centering
%     \includegraphics[width=0.85\linewidth]{certamenes/normal.png}
% \end{figure}


\begin{figure}
    \centering
    \includegraphics[width=0.85\linewidth]{certamenes/normal_1cola.png}
\end{figure}



\end{document}